\documentclass[12pt]{article}

\usepackage{amsmath}
\usepackage{amssymb}
\usepackage{anyfontsize}
\usepackage[toc,page]{appendix}
\usepackage{array}
\usepackage{authblk}
\usepackage{booktabs}
\usepackage{caption}
\usepackage{cellspace}
\usepackage{color}
\usepackage{enumerate}
\usepackage{enumitem}
\usepackage{fancyhdr}
\usepackage{float}
\usepackage{geometry}
\usepackage{graphicx}
\usepackage{listings}
\usepackage{multicol}
\usepackage{newtxtext, newtxmath}
\usepackage{physics} 
\usepackage{subfigure}
\usepackage{tabularx}
\usepackage{wallpaper}

% \CenterWallPaper{.2}{nthu-logo.jpg}

\geometry{
    a4paper,
    total = {170mm, 257mm},
    left = 20mm,
    top = 20mm,
    }

\definecolor{dkgreen}{rgb}{0,0.6,0}
\definecolor{gray}{rgb}{0.5,0.5,0.5}
\definecolor{mauve}{rgb}{0.58,0,0.82}
\lstset{frame = tb,
        language = Python,
        aboveskip = 3mm,
        belowskip = 3mm,
        showstringspaces = False,
        columns = flexible,
        basicstyle = {\small\ttfamily},
        numbers = left,
        numberstyle = \tiny\color{gray},
        keywordstyle = \color{blue},
        commentstyle = \color{dkgreen},
        stringstyle = \color{mauve},
        breaklines = true,
        breakatwhitespace = true,
        tabsize = 4}


%%%%%%%%%%%%%%%%%%%%%%%%%%%%%%%%%%%%%%%%%%%%%%% Again, Don't change anything Above %%%%%%%%%%%%%%%%%%%%%%%%%%%%%%%%%%%%%%%%%%%%%%%


\begin{document}

\pagestyle{fancy}
    \fancyhf{} % clear all header and footer fields
    \renewcommand{\headrulewidth}{0pt} %clear the header line 
    \fancyhead[R]{
        \ifnum\value{page}>1 \includegraphics[width = .23\textwidth]{../figure/nthu-logo.png} 
        \else \fi
        }
    \fancyfoot[C]{\thepage}

% title and info
\title{{\vspace{-2.0cm}
        \normalsize Collider Physics - Final Report}\\
        \textbf{Atmospheric Neutrino Oscillation in Super-K}}
\author{Yuan-Yen Peng}
\affil{\emph{Dept. of Physics, NTHU}\\
        \emph{Hsinchu, Taiwan}}
\date{\today}
\thispagestyle{fancy}
\maketitle

\setlength{\columnsep}{0.03\textwidth}
\begin{multicols}{2}
    
\section{Motivation}
    We comprehend that the government is obligated to upgrade the whole prosperity level of society. Some stages, thereby, ought to be fulfilled; simultaneously, requiring some counterpoise with different hierarchies, inevitably. Some utilitarians may assert that it must quell the low class of people so as to haul up the entire echelon of humankind. We, however, cast doubt on this argument, reckoning that the wipe-out method is too arbitrary and transgressed as a viewpoint of humanity. Ergo, we want to use the numerical model to subvert utilitarianism's implementation.
    
\section{Introduction}
    First, we build a city model and put people into this playground (here, we amend the method, which we presented on the slide sharing part. We recently set people wholly occupied the playground instead of scattering on random positions). Assigned playgrounders with the social rank randomly (instead of setting them according to where they were initially born), and we implement the exchange method (i.e., some rules) to simulate social mobility. After some period, we expect some stable stages, and we will observe the population distribution such as bell distribution, discussing some analogy “physical” quantities. Lastly, apply an “invisible hand” to interrupt the system; in other words, wipe out the low-class people, etc. Afterward, observe the analogical quantities, too.
    First, we build a city model and put people into this playground (here, we amend the method, which we presented on the slide sharing part. We recently set people wholly occupied the playground instead of scattering on random positions). Assigned playgrounders with the social rank randomly (instead of setting them according to where they were initially born), and we implement the exchange method (i.e., some rules) to simulate social mobility. After some period, we expect some stable stages, and we will observe the population distribution such as bell distribution, discussing some analogy “physical” quantities. Lastly, apply an “invisible hand” to interrupt the system; in other words, wipe out the low-class people, etc. Afterward, observe the analogical quantities, too.
    First, we build a city model and put people into this playground (here, we amend the method, which we presented on the slide sharing part. We recently set people wholly occupied the playground instead of scattering on random positions). Assigned playgrounders with the social rank randomly (instead of setting them according to where they were initially born), and we implement the exchange method (i.e., some rules) to simulate social mobility. After some period, we expect some stable stages, and we will observe the population distribution such as bell distribution, discussing some analogy “physical” quantities. Lastly, apply an “invisible hand” to interrupt the system; in other words, wipe out the low-class people, etc. Afterward, observe the analogical quantities, too.

\section{Methodology}
    \subsection{City}
    In reality, the source usage rate depends on distance, constructed a concentric structure, for example, Beijing (Figure), Chicago, etc. Here, our artificially built city will likewise develop around the center of sources; thus, we adopt the concentric zone model to discuss the hierarchy distribution.
    
    % \begin{figure}[H]
    %     \centering 
    %     \includegraphics[width = 0.45\textwidth]{./RK4/0.png}
    %     \includegraphics[width = 0.45\textwidth]{./RK4/2.png} 
    % \end{figure}
    % \begin{figure}[H]
    %     \centering 
    %     \includegraphics[width = 0.45\textwidth]{./RK4/4.png}  
    %     \includegraphics[width = 0.45\textwidth]{./RK4/6.png}
    %     \includegraphics[width = 0.45\textwidth]{./RK4/8.png}
    %     \includegraphics[width = 0.45\textwidth]{./RK4/10.png} 
    %     \caption{This the normal cloud with RK4, from $t=0$ to $t=10$ and $dt=0.01$.}
    %     \label{RK4_cloud}
    % \end{figure}

    \subsection{Leapfrog method}
    
    \subsection{Energy comparison}
    
\end{multicols}


\end{document}